\section{Background} \label{sec:background}

\subsection{Discrete Systems and Cotyledons}
In Wymore's framework, a discrete system (denoted $\mathrm{DSYSTEM}$ in \citet{Wymore1993}) is a deterministic Moore-type state machine defined by a quintuple $Z=(S,I,O,\delta,\rho)$ consisting of a nonempty state set $S$, an input alphabet $I$, an output alphabet $O$, a next-state transition function $\delta:S\times I\to S$, and a state readout function $\rho:S\to O$. Here ``discrete'' refers to the discrete time scale $T=\mathbb{N}$; the state and port sets $S$, $I$, and $O$ may be either finite (as in physical digital hardware) or countably infinite (as in idealized mathematical specifications with, e.g., unbounded counters). Given an initial state $s_0\in S$ and an input trajectory $f:T\to I$, iterating $\delta$ from $s_0$ generates a unique state trajectory $g(t)$ (Wymore's $\mathrm{STZ}(Z,s_0,f)(t)$) and a corresponding boundary output trajectory $y(t)=\rho(g(t))$ (Wymore's $\mathrm{OTZ}(Z,s_0,f)(t)$). Because $\rho$ is a function of the current state $g(t)$ alone, every Wymorian discrete system obeys Moore causality: an external input $f(t)$ arriving at tick $t$ first affects the next state $g(t+1)=\delta(g(t),f(t))$ and can therefore influence the boundary output no earlier than tick $t+1$.

T3SD organizes a system design problem around three mathematical spaces, called \emph{cotyledons} after their leaf-shaped depiction in Wymore's diagrams \citep{Wymore1993}:
\begin{enumerate}
    \item \textbf{Functionality Cotyledon $\mathrm{CTL}(IOR)$.} An input/output requirement is a 6-tuple $IOR=(OLR,IR,ITR,OR,OTR,ER)$, where $OLR$ specifies the operational length (defining the required time scale $T_{IOR}\subseteq\mathbb{N}$), $IR$ and $OR$ are the boundary input and output alphabets, $ITR\subseteq IR^T$ and $OTR\subseteq OR^T$ are the sets of admissible input and output trajectories, and $ER\subseteq ITR\times OTR$ is the eligibility relation pairing each admissible input trajectory $f\in ITR$ with its set of acceptable output trajectories $ER(f)\subseteq OTR$. A \emph{functional system design} $D_f=(Z_f,s_0,T_f)\in\mathrm{CTL}(IOR)$ (the Functionality Space Requirement, $\mathrm{FSR}(IOR)$) packages a discrete system $Z_f$ on boundary alphabets $(IR,OR)$, a designated initial state $s_0$, and a nonempty observation horizon $T_f\subseteq T_{IOR}$ such that for every $f\in ITR$ (and any input extension agreeing with $f$ on $T_{IOR}$), the output trajectory generated by $Z_f$ from $s_0$ agrees on $T_f$ with some eligible trajectory $g\in ER(f)$. Wymore's shows that $\mathrm{CTL}(IOR)$ is either empty (when $IOR$ violates Moore causality) or infinite, because once a single realizing state machine exists, any state re-encoding or state-space elaboration yields another distinct member of $\mathrm{CTL}(IOR)$ (see Theorem~6.64 in \citep{Wymore1993}).
    \item \textbf{Buildability Cotyledon $\mathrm{CTL}(TYR)$.} A technology $TYR$ is a nonempty catalog of available discrete component systems (such as hardware registers, sensors, clocks, or software modules). A System Coupling Recipe $SCR=(VSCR,CSCR)$ specifies a vector $VSCR$ of component systems together with an internal port-to-port feedback and feedforward wiring relation $CSCR$. When every component in $VSCR$ is selected from the catalog $TYR$, synchronous composition of the wired components produces a unique coupled \emph{resultant} discrete system $Z_b=\mathrm{RSY}(SCR)$ whose external boundary inputs and outputs are the unconnected input and output ports of the recipe (denoted $\mathrm{UISCR}(SCR)$ and $\mathrm{UOSCR}(SCR)$ in \citet{Wymore1993}). The buildability cotyledon $\mathrm{CTL}(TYR)$ (the Buildability Space Requirement, $\mathrm{BSR}(TYR)$) consists of all such buildable designs $D_b=(SCR,Z_b)$ constructible from $TYR$.
    \item \textbf{Implementability Cotyledon $\mathrm{CTL}(IOR,TYR)$.} Informally, we may uderstand the implementability cotyledon as the intersection between functionality and buildability, i.e., those systems that both satisfy the input/output requirement and are constructible from the available technology. Formally, because functional designs specify abstract boundary behavior whereas buildable designs specify coupled component architectures with potentially different state spaces and port structures, $\mathrm{CTL}(IOR,TYR)$ (the Implementability Space Requirement, $\mathrm{ISR}(IOR,TYR)$) is defined as the set of pairs $(D_f,D_b)\in\mathrm{CTL}(IOR)\times\mathrm{CTL}(TYR)$ linked by Wymore's \emph{implementation} relation ($\mathrm{IMPSY}$, Definition~5.71) \citep{Wymore1993}.
\end{enumerate}
Throughout this paper, when the initial state $s_0$, observation horizon $T$, and coupling recipe $SCR$ are clear from context, we identify a functional design $D_f=(Z_f,s_0,T)$ with its underlying system $Z_f\in\mathrm{CTL}(IOR)$ and a buildable design $D_b=(SCR,Z_b)$ with its coupled resultant $Z_b=\mathrm{RSY}(SCR)\in\mathrm{CTL}(TYR)$, writing implementable pairs in $\mathrm{CTL}(IOR,TYR)$ as $(D_f,Z_b)$.

\subsection{System Modes, Homomorphisms, and Implementation}
\label{sec:background-modes-hom}
We would rarely expect a buildable resultant $Z_b=\mathrm{RSY}(SCR)$ to satisfy an abstract functional model $Z_f\in\mathrm{CTL}(IOR)$ by direct state-space equality, as at the very least, the two systems will almost certainly have different state spaces.  In practice, a coupled architecture contains implementation details that do not appear in an idealized functional specification: multi-phase internal clocks, pipeline latches, post-reset initialization states, diagnostic registers, or unused port encodings. Conversely, an idealized functional model may be stated over an infinite state space (such as an unbounded counter), whereas any physical machine operates within a bounded hardware capacity.

Wymore introduces two concepts to relate buildable resultants and functional models:
\begin{itemize}
    \item \textbf{System Mode ($\mathrm{SYSMO}$, see Definition~5.6 in \citet{Wymore1993}).} A discrete system $Z_m=(S_m,I_m,O_m,\delta_m,\rho_m)$ is a \emph{system mode} of a buildable system $Z_b=(S_b,I_b,O_b,\delta_b,\rho_b)$, written $Z_m\in\mathrm{SYSMO}(Z_b)$, if $Z_m$ embeds into $Z_b$ via injective maps on states, inputs, and outputs such that each step of $Z_m$ corresponds to a well-defined trajectory segment of $Z_b$ within an invariant operational regime (for example, restricting $Z_b$ to its reachable post-reset states, an active clock phase, or a bounded operating envelope). Every system is trivially a mode of itself via its primary self-mode.
    \item \textbf{System Homomorphism ($\mathrm{HIMSY}$, see Definition~4.3 in \citet{Wymore1993}).} A functional system $Z_f=(S_f,I_f,O_f,\delta_f,\rho_f)$ is a \emph{homomorphic image} of $Z_m$ (written $Z_f\in\mathrm{HIMSY}(Z_m)$) if there exist \emph{surjective} maps $h_S:S_m\twoheadrightarrow S_f$, $h_I:I_m\twoheadrightarrow I_f$, and $h_O:O_m\twoheadrightarrow O_f$ (denoted $\mathrm{HS}$, $\mathrm{HI}$, and $\mathrm{HO}$ in \citet{Wymore1993}) that commute with state transitions and readouts at every state $x\in S_m$ and input $p\in I_m$:
    \begin{equation}
        h_S(\delta_m(x,p)) = \delta_f(h_S(x),h_I(p)), \qquad h_O(\rho_m(x)) = \rho_f(h_S(x)).
        \label{eq:wymore-hom-def}
    \end{equation}
\end{itemize}
By Definition~5.71 in \citet{Wymore1993}, a buildable resultant $Z_b$ \emph{implements} a functional system $Z_f$ (written $Z_f\in\mathrm{IMPSY}(Z_b)$) if and only if $Z_f$ is a homomorphic image of at least one system mode $Z_m\in\mathrm{SYSMO}(Z_b)$. Wymore's Theorem~6.58 further establishes that when $Z_m$ and $Z_f$ share the boundary alphabets $(IR,OR)$ with identity port maps ($h_I=\mathrm{id}_{IR}$, $h_O=\mathrm{id}_{OR}$), any homomorphism $h_S:S_m\twoheadrightarrow S_f$ preserves $IOR$ satisfaction across corresponding initial states:
\begin{equation}
    \mathrm{SatisfiesIOR}(Z_m,IOR,s_0,T) \Longleftrightarrow \mathrm{SatisfiesIOR}(Z_f,IOR,h_S(s_0),T).
    \label{eq:thm-658-preview}
\end{equation}

\miniheading{Notation}
Table~\ref{tab:notation} summarizes Wymore's textbook nomenclature \citep{Wymore1993} alongside the standard mathematical notation used throughout this paper.

\begin{table*}[t]
\caption{Correspondence between Wymore's T3SD terminology \citep{Wymore1993} and the mathematical notation used in this paper.}
\label{tab:notation}
\renewcommand{\arraystretch}{1.25}
\footnotesize
\begin{tabularx}{\textwidth}{@{}A{0.25\textwidth} A{0.35\textwidth} >{\raggedright\arraybackslash}X@{}}
    \toprule
    \textbf{Concept} & \textbf{Wymore (T3SD)} & \textbf{Notation in This Paper} \\
    \midrule
    Discrete system 5-tuple & $\mathrm{DSYSTEM}\ (SZ, IZ, OZ, \mathrm{NZ}, \mathrm{RZ})$ & $Z = (S, I, O, \delta, \rho)$ \\
    State / output trajectories & $\mathrm{STZ}(Z,s_0,f)(t),\ \mathrm{OTZ}(Z,s_0,f)(t)$ & $g(t),\ y(t)=\rho(g(t))$ \\
    Input/Output Requirement & $\mathrm{IOR} = (\mathrm{OLR},\mathrm{IR},\mathrm{ITR},\mathrm{OR},\mathrm{OTR},\mathrm{ER})$ & $IOR = (OLR, IR, ITR, OR, OTR, ER)$ \\
    Functionality Cotyledon & $\mathrm{CTL}(\mathrm{IOR}) = \mathrm{FSR}(\mathrm{IOR})$ & $\mathrm{CTL}(IOR) = \mathrm{FSR}(IOR)$ \\
    Buildability Cotyledon & $\mathrm{CTL}(\mathrm{TYR}) = \mathrm{BSR}(\mathrm{TYR})$ & $\mathrm{CTL}(TYR) = \mathrm{BSR}(TYR)$ \\
    Implementability Cotyledon & $\mathrm{CTL}(\mathrm{IOR}, \mathrm{TYR}) = \mathrm{ISR}(\mathrm{IOR},\mathrm{TYR})$ & $\mathrm{CTL}(IOR, TYR) = \mathrm{ISR}(IOR, TYR)$ \\
    Coupling recipe \& resultant & $\mathrm{SCR} = (\mathrm{VSCR}, \mathrm{CSCR}),\ \mathrm{RSY}(\mathrm{SCR})$ & $SCR = (VSCR, CSCR),\ Z_b = \mathrm{RSY}(SCR)$ \\
    System modes of $Z_b$ & $\mathrm{SYSMO}(Z_b)$ & $Z_m \in \mathrm{SYSMO}(Z_b)$ \\
    Homomorphic image / maps & $Z_f \in \mathrm{HIMSY}(Z_m)$ via $(\mathrm{HS},\mathrm{HI},\mathrm{HO})$ & $h = (h_S, h_I, h_O) : Z_m \twoheadrightarrow Z_f$ \\
    Implementation relation & $Z_f \in \mathrm{IMPSY}(Z_b)$ & $\exists Z_m \in \mathrm{SYSMO}(Z_b),\ h : Z_m \twoheadrightarrow Z_f$ \\
    Temporal specifications & --- & $Z_m\models_{IO}\Phi_{IOR}$ (boundary),\ $Z_m\models_{\mathrm{dyn}}\Phi_{\mathrm{dyn}}(Z_{\mathrm{ref}})$ (dynamics) \\
    \bottomrule
\end{tabularx}
\end{table*}

\subsection{Temporal Logics and Model Checking}
\label{sec:background-ltl}
Linear Temporal Logic (LTL) and its first-order extensions provide a formal language for expressing properties of time-indexed trajectories \citep{Pnueli1977,MannaPnueli1995,BaierKatoen2008,VardiWolper1986,Thomas1990}. Rather than specifying desired behavior by constructing a concrete state set and transition table, a temporal formula characterizes the admissible infinite or finite traces directly using Boolean connectives ($\wedge,\vee,\neg,\rightarrow$) and temporal modalities such as $G\,\varphi$ (``always $\varphi$'') and $X\,\varphi$ (``at the next tick, $\varphi$''), summarized in Table~\ref{tab:ltl-operators}.

\begin{table}[t]
\caption{Core temporal logic operators used to express boundary requirements $\Phi_{IOR}$ and compiled dynamics $\Phi_{\mathrm{dyn}}$.}
\label{tab:ltl-operators}
\renewcommand{\arraystretch}{1.2}
\footnotesize
\setlength{\tabcolsep}{4pt}
\begin{tabularx}{\columnwidth}{@{}l l >{\raggedright\arraybackslash}p{0.28\columnwidth} >{\raggedright\arraybackslash}X@{}}
    \toprule
    \textbf{Op.} & \textbf{Name} & \textbf{Role in T3SD} & \textbf{Semantics at Tick $t$} \\
    \midrule
    $G\,\varphi$ & always & Global invariant & $\varphi$ holds at all $\tau \ge t$ \\
    $X\,\varphi$ & next & 1-step delay & $\varphi$ holds at $t+1$ \\
    $\varphi \rightarrow \psi$ & implies & Guarded step & If $\varphi$ at $t$, then $\psi$ at $t$ \\
    $\varphi \wedge \psi$ & and & Joint guard & Both $\varphi$ and $\psi$ hold at $t$ \\
    $F\,\varphi$ & eventually & Attainment & $\varphi$ holds at some $\tau \ge t$ \\
    \bottomrule
\end{tabularx}
\end{table}

When a discrete system $Z_m$ has finite state, input, and output sets, its execution defines a finite Kripke structure whose paths are infinite sequences of $(s(t),i(t),y(t))$ triples. Deciding whether every run of a finite-state system satisfies an LTL formula ($Z_m\models\varphi$) is PSPACE-complete \citep{SistlaClarke1985} and supported by established explicit-state, symbolic binary decision diagram (BDD), and SAT/SMT-based model checkers \citep{BurchClarkeMcMillanDill1992,BiereCimattiClarkeZhu1999,BaierKatoen2008,deMouraZ32008}.

\subsection{Related Work}
\label{sec:related-work}
Several lines of research across systems engineering, formal verification, and computational complexity inform the approach developed in this paper.

Within the MBSE literature, recent work has revisited Wymore's T3SD to provide a rigorous systems-theoretic foundation for modern modeling practices. Salado and Wach \citep{salado2019wymore,salado2019mbse} demonstrated how Wymorian input/output requirements can ground True Model-Based Requirements (TMBR) in SysML without conflating problem-space boundary requirements with internal solution architectures, while Wach et al.\ \citep{wach2021devs} connected Wymore's discrete systems to Discrete Event System Specification (DEVS) \citep{ZeiglerPraehoferKim2000}. Mainstream SysML/UML practice, however, lacks a native mathematical notion of homomorphism between problem-space boundary requirements and coupled buildable resultants, so architectural verification often remains descriptive or else exports the entire diagram into an external semantic substrate such as Petri nets or transition systems \citep{Jensen1997,ZeiglerPraehoferKim2000}.  T3SD-oriented modernization efforts preserve the separation between the functionality and buildability cotyledons; however, they still rely on Wymore's classical constructive $\mathrm{IMPSY}$ definitions when certifying that a coupled architecture implements a requirement. The present paper retains Wymore's coupling algebra $\mathrm{RSY}$ and system modes $\mathrm{SYSMO}$ and automates where possible the implementability check through temporal specifications $\Phi_{IOR}$ and $\Phi_{\mathrm{dyn}}$.

Contract-based design and interface theories provide a complementary vocabulary for boundary requirements.  Assume/guarantee contracts and interface automata typically partition concerns into environmental assumptions and system guarantees \citep{BenvenisteEtAl2018,deAlfaroHenzinger2001}. A Wymorian $IOR$, by contrast, packages admissible input trajectories $ITR$ together with an eligibility relation $ER\subseteq ITR\times OTR$ that pairs each admissible input with its acceptable outputs. Expressing that eligibility relation as a boundary temporal specification $\Phi_{IOR}$ makes $IOR$ satisfaction tool-checkable in the spirit of contract-style boundary verification, while retaining Wymore's Moore causality at the readout boundary; we do not identify $IOR$s with assume/guarantee contracts, but rather note that $\Phi_{IOR}$ supplies a temporally executable account of the same trajectory-level intent that contracts address at interfaces.

In formal methods and concurrency theory, temporal logic safety fragments have long been used to characterize step-local transition invariants \citep{Sistla1994,MannaPnueli1995}, and characteristic-formula constructions encode operational state machines into modal or temporal logic formulas up to simulation or bisimulation equivalence \citep{Milner1989,GrafSifakis1986,SteffenIngolfsdottir1994,AcetoEtAl2012}. Those constructions typically target action-labeled transition systems and often rely on the modal calculi to characterize branching or weak behavioral equivalences. Algebraically, Wymore's surjective system homomorphism corresponds to a bounded morphism on deterministic Moore coalgebras \citep{Rutten2000,BlackburnDeRijkeVenema2001}. T3SD systems, however, are deterministic Moore machines with synchronous port-based coupling ($\mathrm{RSY}$) and operational modes ($\mathrm{SYSMO}$) that preserve step-local autonomous behavior. Consequently, the dynamics fragment $\Phi_{\mathrm{dyn}}$ is intentionally a step-local $G$/$X$ safety fragment that mirrors Wymore's static homomorphism axioms (see Definition~4.3 in \citet{Wymore1993}), rather than a full characteristic formula for branching bisimulation. However, prior characteristic-formula results do not, out of the box, connect trajectory-based $IOR$s, system modes, and coupling recipes within the tricotyledon framework.  This is a gap that this paper addresses.

Finally, classical refinement theories and combinatorial complexity delimit what lock-step T3SD verification can and cannot claim. Refinement mappings and hierarchical I/O automata permit internal stuttering or silent steps between corresponding observable transitions \citep{AbadiLamport1991,LynchTuttle1987,Lamport2002}. Wymore's implementation relation instead equates one tick of the buildable mode with one tick of the functional design: intermediate execution artifacts are handled by first isolating an operational system mode $Z_m\in\mathrm{SYSMO}(Z_b)$ and then exhibiting a surjective step-synchronous homomorphism onto $Z_f$. The compiled fragment $\Phi_{\mathrm{dyn}}$, which retains the next-step operator $X$, is therefore the exact temporal counterpart of classical T3SD.  Multi-rate and stuttering relaxations are not considered in this work. Deciding whether a surjective homomorphism exists between finite relational structures or transition systems is NP-hard in general \citep{HellNesetril1990,BodirskyKaraMartin2012}, and maintaining an explicit state-space projection $h_S:S_m\twoheadrightarrow S_f$ is fragile under architectural revision. By translating $IOR$s and reference dynamics into temporal specifications evaluated over finite buildable system modes, our work replaces manual homomorphism construction with decidable LTL and SAT-based checking \citep{SistlaClarke1985,BaierKatoen2008}.
